\documentclass[10pt]{article}
\usepackage[T1]{fontenc}
\usepackage[utf8]{inputenc}
\usepackage[french]{babel}
\usepackage{amsmath,amssymb,mathrsfs,mathpazo}
\usepackage{geometry}
\geometry{a4paper,margin=2cm}
\usepackage{graphicx,array,tabularx,tabularray,colortbl,multirow,multicol,booktabs,enumitem,pifont,ulem,fancybox,diagbox,cancel}
\usepackage{pstricks,pstricks-add,pst-plot,pst-tree,pst-node,pst-text,pst-func,pst-math,pst-bspline,pst-3dplot,pst-eucl,pst-solides3d}
\let\pscalculate\undefined
\usepackage{pst-bezier}
\usepackage{tikz,pgfplots,tkz-tab}
\usetikzlibrary{arrows,arrows.meta,patterns,calc,positioning,trees,decorations.markings,shapes.geometric,shapes.misc,shapes.arrows,intersections,decorations.pathmorphing,decorations.pathreplacing,backgrounds,shadows,babel}
\pgfplotsset{compat=1.18}
\usepgfplotslibrary{fillbetween}
\usepackage[np]{numprint}
\usepackage{listings,makecell,pas-tableur,twoopt}
\let\barre\undefined
\usepackage[pstricks]{tablvar}
\definecolor{darkgreen}{RGB}{0,100,0}
\providecommand{\red}{\color{red}}
\providecommand{\blue}{\color{blue}}
\providecommand{\green}{\color{green}}
\providecommand{\black}{\color{black}}
\providecommand{\white}{\color{white}}
\providecommand{\gray}{\color{gray}}
\providecommand{\cyan}{\color{cyan}}
\providecommand{\magenta}{\color{magenta}}
\providecommand{\yellow}{\color{yellow}}
\providecommand{\violet}{\color{violet}}
\providecommand{\orange}{\color{orange}}
\providecommand{\euro}{\texteuro}
\providecommand{\EUR}{\texteuro}
\providecommand{\ieme}{\textsuperscript{e}}
\providecommand{\degres}{\ensuremath{^\circ}}
\providecommand{\degree}{\ensuremath{^\circ}}
\providecommand{\vv}[1]{\overrightarrow{#1}}
\providecommand{\Coord}[1]{#1}
\providecommand{\decosix}{}
\providecommand{\starredbullet}{\ensuremath{\bullet}}
\pagestyle{empty}
\setlength{\parindent}{0pt}
\usepackage{adjustbox}\DeclareUnicodeCharacter{2236}{:}\DeclareUnicodeCharacter{2190}{\ensuremath{\leftarrow}}
\providecommand{\R}{\mathbb{R}}
\providecommand{\N}{\mathbb{N}}
\providecommand{\Z}{\mathbb{Z}}
\providecommand{\Q}{\mathbb{Q}}
\providecommand{\D}{\mathbb{D}}
\providecommand{\C}{\mathbb{C}}
\providecommand{\e}{{\mathrm{e}}}
\providecommand{\vect}[1]{\overrightarrow{#1}}
\providecommand{\vectt}[1]{\overrightarrow{\mathrm{#1}}}
\providecommand{\vvt}[1]{\overrightarrow{\mathrm{#1}}}
\providecommand{\np}[2][]{{#2~#1}}
\providecommand{\vv}[1]{{\overrightarrow{#1}}}
\renewcommand{\d}{{\mathrm{d}}}
\providecommand{\degres}{{^\circ}}
\providecommand{\degre}{{^\circ}}
\providecommand{\blue}{\color{blue}}
\providecommand{\red}{\color{red}}
\providecommand{\enskip}{\quad}
\providecommand{\llbracket}{\lbrack\!\lbrack}
\providecommand{\rrbracket}{\rbrack\!\rbrack}
\providecommand{\euro}{{\text{€}}}
\providecommand{\no}{{\text{n°}}}
\providecommand{\interval}[2]{{[#1;#2]}}
\providecommand{\intervFO}[2]{{[#1;#2[}}
\providecommand{\dotfill}{\ldots}

\providecommand{\ds}{\displaystyle}
\providecommand{\textding}[1]{#1}
\providecommand{\Oij}{\(\mathrm{O};\vec{\imath},\vec{\jmath}\)}
\providecommand{\Oijk}{\(\mathrm{O};\vec{\imath},\vec{\jmath},\vec{k}\)}
\providecommand{\Ouv}{\(\mathrm{O};\vec{u},\vec{v}\)}
\providecommand{\intervFF}[2]{\ensuremath{[#1;#2]}}
\renewcommand{\euro}{\ifmmode\text{\texteuro}\else\texteuro\fi}
\DeclareUnicodeCharacter{00B0}{\ifmmode^{\circ}\else\textdegree\fi}
\DeclareUnicodeCharacter{20AC}{\ifmmode\text{\texteuro}\else\texteuro\fi}
% Énoncés : Éducation nationale. Transcription et corrigés : APMEP, Denis Vergès et contributeurs. Découpage et classement : CoopMaths, Sébastien Lozano, Éric Elter et contributeurs. Adaptation de la présentation : Maths Explorer.
% CC BY-SA : https://coopmaths.fr/www/about/
% https://coopmaths.fr/alea/static/eam/2025/tex/eam_specialite_2025_06_sujet0v2_2_cor.tex
\begin{document}
\noindent\textit{APMEP / CoopMaths. CC BY-SA. Présentation adaptée par Maths Explorer.}\par\medskip
\medskip

En 2020, une ville comptait 10 000 habitants

On modélise l'évolution du nombre d'habitants de cette ville par la suite $(u_n)$ définie ainsi : 
\[\left\{\begin{array}{l c l}
u_0&=&\np{10000}\\
u_{n+1}&=&1,08u_n - 300,\quad n \in \N 
\end{array}\right.\]

où $u_n$ représente le nombre d'habitants pour l'année $2020 + n$.



\begin{enumerate}
\item
$u_1$  représente le nombre d'habitants pour l'année $2020 + 1$ soit le nombre d'habitants en 2021.

$u_{1}=1,08 u_0 - 300 = 1,08\times \np{10000}-300 = \np{10800}-300=\np{10500}$

\item On considère la suite $(v_n)$ définie pour tout entier naturel $n$ par $v_n = u_n  - \np{3750}$. 
	\begin{enumerate}
		\item 
$v_0=u_0 - 	\np{3750} = \np{10000} - \np{3750} = \np{6250}$
		
		\item 
Pour tout entier naturel $n$, on a $v_n=u_n-\np{3750}$ donc $u_n=v_n+\np{3750}$.

$\begin{aligned}
v_{n+1}
& =u_{n+1}-\np{3750}
 = 1,08u_n -300 -\np{3750}
= 1,08\left ( v_n +\np{3750} \right ) -300 - \np{3750}\\
& =1,08 v_n + \np{4050} - \np{4050}
= 1,08 v_n
\end{aligned}$

		\item
Donc la suite $(v_n)$ est géométrique de raison $q=1,08$ et de premier terme $v_0 = \np{6250}$.

		\item On en déduit que, pour tout entier naturel $n$, on a 
$v_n=v_0\times q^n = \np{6250}\times 1,08^n$.		
		
		\item 
Pour tout $n$ on a $v_n= \np{6250}\times 1,08^n$ et $u_n=v_n+\np{3750}$ donc on peut en déduire que $u_n = \np{6250}\times 1,08^n +\np{3750}$.
	\end{enumerate}
	
\item ~

\begin{minipage}{0.66\linewidth}
Le tableau ci-contre, extrait d'une feuille automatisée de calcul, a été obtenu par recopie vers le bas après avoir saisi la formule suivante dans la cellule B2 :

\begin{center}
\fbox{\texttt{= 6250*1,08\^{}A2 + 3750}}
\end{center}

La municipalité envisage d'ouvrir une nouvelle école
 maternelle dès que la population atteindra \np{19000} habitants.\\

\end{minipage}
\hfill
\begin{minipage}{0.3\linewidth}
$\begin{array}{|>{\cellcolor{lightgray}}c|c|c|}\hline
\rowcolor{lightgray}&\text{A}&\text{B}\\ \hline
1&\text{n}& \text{Un}\\ \hline
2&0&\np{10000}\\ \hline
3&1&\np{10500}\\ \hline
\cdots & \cdots & \cdots \\ \hline
13&11&\np{18322,74373}\\ \hline
14&12&\np{19488,56323}\\ \hline
15&13&\np{20747,64829}\\ \hline
\cdots & \cdots & \cdots \\ \hline
\end{array}$
\end{minipage}

\medskip

La population dépasse $\np{19000}$ habitants pour $n=12$ soit en 2032.

La construction d'un tel établissement nécessite deux ans; il faut donc commencer les travaux en 2030.

\end{enumerate}

\bigskip

\end{document}
